\subsection{Листинг численного метода}

Для реализации блок-схем воспользуемся языком программирования C++.

\begin{lstlisting}[caption=Прямой шаг метода Гаусса.]
bool triangulateMatrix(const int n, vector<vector<double>>& matrix) {
  bool isSingular = false;
  for (size_t i = 0; i < n - 1; i++) {
    size_t idx = n;
    for (size_t j = i; j < n; j++) {
      if (matrix[j][i] != 0) {
        idx = j;
        break;
      }
    }
    if (idx == n) {
      isSingular = true;
      continue;
    }
    // apply P operator
    if (idx != i) {
      vector<double> temp_row = matrix[idx];
      matrix[idx] = matrix[i];
      matrix[i] = temp_row;
    }
    // apply M operator
    for (size_t j = i + 1; j < n; j++) {
      double factor = matrix[j][i] / matrix[i][i];
      for (size_t k = i; k <= n; k++) {
        matrix[j][k] -= factor * matrix[i][k];
      }
    }
  }
  if (isSingular) {
    errorMessage = "The system has no single solution.";
    isSolutionExists = !hasContradiction(n, matrix);
    return false;
  }
  return true;
}
\end{lstlisting}

\begin{lstlisting}[caption=Проверка системы на совместность.]
bool hasContradiction(const int n, vector<vector<double>>& matrix) {
  vector<double> zeroes(n, 0);
  for (size_t i = 0; i < n; i++) {
    vector<double> curr_row = matrix[i];
    curr_row.pop_back();
    if (curr_row == zeroes && matrix[i][n] != 0) {
      return true;
    }
  }
  return false;
}
\end{lstlisting}

\begin{lstlisting}[caption=Обратный шаг метода Гаусса.]
vector<double> solveByGauss(int n, vector<vector<double>> matrix) {
  vector<double> answer(2 * n, 0);
  vector<vector<double>> matrix_triangle = matrix;
  if (!triangulateMatrix(n, matrix_triangle)) {
    return answer;
  }
  for (size_t i = 0; i < n; i++) {
    const size_t idx = n - i - 1;
    for (size_t j = idx + 1; j < n; j++) {
      answer[idx] -= matrix_triangle[idx][j] * answer[j]; 
    }
    answer[idx] += matrix_triangle[idx][n];
    answer[idx] /= matrix_triangle[idx][idx];
  }
  // calculate residuals
  for (size_t i = 0; i < n; i++) {
    for (size_t j = 0; j < n; j++) {
      answer[n + i] += matrix[i][j] * answer[j];
    }
    answer[n + i] -= matrix[i][n];
  }
  return answer;
}
\end{lstlisting}